\section{All tensor powers of the flip--cup algebra in wire dimension four}
\label{sec:full-flip-cup}

The finite two-copy calculation leaves open whether different
coefficients at larger copy numbers could improve the exponent.
An invariant-support argument resolves that question for the entire
flip--cup algebra in wire dimension four. The bound applies to arbitrary
coefficients in both summands, including coefficients that change with
the number of copies. It is not a statement about unrestricted
partial-transpose decompositions.

The algebra generated by the identity, flip and cup is the familiar
algebra of orthogonally invariant operators; its three projectors and
their partial transposes are described in
\cite[Section~VI.B, pp.~5--6]{ADVW2002}. That work concerns positive
states and entanglement quantities. We impose no positivity condition
and prove the rank bound below directly.

\begin{lemma}[Invariant supports and product packings]
\label{lem:invariant-support-packing}
Let $T$ be a linear involution on $\kk^J$, where $J$ is finite,
and suppose $T\mathbf1=\mathbf1$. Assign positive real weights
$w_x$ to $x\in J$. Suppose linear functionals $\ell_j$ satisfy
\[
 \ell_jT=\ell_j,\qquad \ell_j(\mathbf1)=1,
 \qquad \operatorname{supp}\ell_j\subseteq E_j\subseteq J.
\]
Let $\lambda_j\geq0$ be real numbers such that
\[
 \sum_{j:x\in E_j}\lambda_j\leq w_x\qquad(x\in J),
 \qquad L=\sum_j\lambda_j.
\]
If $m,n\in\kk^{J^t}$ satisfy
$m+T^{\otimes t}n=\mathbf1$, then
\[
 \sum_{x\in\operatorname{supp}m}\prod_{a=1}^t w_{x_a}
 +\sum_{x\in\operatorname{supp}n}\prod_{a=1}^t w_{x_a}
 \geq L^t.
\]
If $\operatorname{char}\kk\ne2$ and
$a+T^{\otimes t}a=\mathbf1$, the same lower bound holds for the
single weighted support of $a$. The conclusion also holds for a
product of different finite systems, with the product of their
packing totals on the right.
\end{lemma}
\begin{proof}
For every tuple $\mathbf j=(j_1,\ldots,j_t)$, the tensor product
$\ell_{j_1}\otimes\cdots\otimes\ell_{j_t}$ is invariant under
$T^{\otimes t}$ and takes value one on $\mathbf1$.
Thus its values on $m,n$ sum to one. At least one of the two vectors
has a nonzero coordinate in
$E_{j_1}\times\cdots\times E_{j_t}$.
In the self-complementary case its value on $a$ is $1/2$, so the
same product set meets the support of $a$.

Give this product set weight $\prod_a\lambda_{j_a}$.
Their total weight is $L^t$. Each support coordinate $x$ can account
for at most
\[
 \sum_{\mathbf j:x_a\in E_{j_a}\ \forall a}
       \prod_a\lambda_{j_a}
 =\prod_a\left(\sum_{j:x_a\in E_j}\lambda_j\right)
 \leq\prod_a w_{x_a}
\]
of that weight. Summing over the support, or the two supports,
proves the inequalities. The argument uses no tensorization claim
about minimum rank. It works unchanged with different systems
in the factors.
\end{proof}

\begin{theorem}[The full algebra in wire dimension four stays above the square exponent]
\label{prop:full-flip-cup-four}
Let $\kk$ have characteristic other than two. On
$\kk^4\otimes\kk^4$, let $F$ be the flip and
$\Omega=vv^{\mathsf T}$, where $v=\sum_{i=1}^4e_i\otimes e_i$.
Put $\mathcal B_4=\operatorname{span}\{I,F,\Omega\}$, and let
$\Gamma$ be partial transpose in the second factor, extended
copywise. For every $k\geq0$, every
$A\in\mathcal B_4^{\otimes k}$ satisfying $A+\Gamma A=I$ obeys
\begin{equation}\label{eq:full-flip-cup-four-bound}
 \operatorname{rank}A\geq
 52^{\lfloor k/2\rfloor}7^{k\bmod2}.
\end{equation}
For every $M,N$ in the same algebra with $I=M+\Gamma N$,
their total rank obeys the same bound.

Consequently every sequence of such decompositions has lower
asymptotic rank rate at least $\sqrt{52}$ and exponent, in wire
dimension $4^k$, at least
\[
 \tfrac12\log_4 52=1.4251099\ldots>\sqrt2.
\]
\end{theorem}
\begin{proof}
The simultaneous invariant spaces are the cup line, the symmetric
trace-zero space, and the antisymmetric space. Denote their types
by $0,+,-$; their dimensions are respectively $1,9,6$.
This is a direct sum over every field under consideration: the
eigenvalues of $F$ are $1,-1$, and $\Omega$ acts by the nonzero
scalar four on the cup line and by zero on its complement.
An element of $\mathcal B_4$ is scalar on each of these spaces.

Partial transpose fixes $I$ and exchanges $F,\Omega$.
If an operator has eigenvalues $(a_0,a_+,a_-)$, direct substitution
in $\alpha I+\beta F+\gamma\Omega$ shows that its partial
transpose has eigenvalue vector $Ta$, where
\begin{equation}\label{eq:flip-cup-four-transform}
 T=\frac14\begin{pmatrix}
 1&9&-6\\1&1&2\\-1&3&2
 \end{pmatrix}.
\end{equation}
In $k$ copies, eigenvalue coordinates are indexed by $\{0,+,-\}^k$,
their block dimensions are the products of $w=(1,9,6)$, and the
transformation is $T^{\otimes k}$. Matrix rank is exactly the
weighted support of the eigenvalue vector.

For two copies use the coordinate order
\[
 00,\ 0+,\ 0-,\ +0,\ ++,\ +-,\ -0,\ -+,\ --.
\]
The following five rows define the functionals $8\ell_1,\ldots,8\ell_5$:
\begin{equation}\label{eq:flip-cup-four-certificate}
 \begin{pmatrix}
 -1&0&0&0&-3&6&0&6&0\\
 0&1&1&0&3&3&0&0&0\\
 0&0&0&1&3&0&1&3&0\\
 0&0&0&0&2&2&0&2&2\\
 0&1&0&1&2&2&0&2&0
 \end{pmatrix}.
\end{equation}
Every row sums to eight. Multiplying each row by $T\otimes T$
leaves it unchanged, as direct multiplication of the displayed
integer matrices verifies. These identities only divide by powers
of two, so they hold in every characteristic in the theorem.
Take $E_j$ to be the displayed rational support of row $j$, even
if some of its entries vanish after reduction, and assign weights
\[
 (\lambda_1,\ldots,\lambda_5)=(1,6,6,36,3).
\]
Their total is $52$. Their coordinate loads are
\[
 (1,9,6,9,52,46,6,46,36)
 \leq(1,9,6,9,81,54,6,54,36)=w\otimes w.
\]
Lemma~\ref{lem:invariant-support-packing} applied to $t$ such
two-copy systems proves the factor $52^t$.

For a remaining single copy, use
$\ell_0=(1,3,0)/4$ and $\ell_1=(0,2,2)/4$.
Both are invariant under $T$ and take value one on $\mathbf1$.
Assign their supports weights one and six. The coordinate loads
are $(1,7,6)\leq(1,9,6)$, with total seven. The mixed product
version of the lemma proves \eqref{eq:full-flip-cup-four-bound}
for odd $k$ as well. For $k=0$, the self-complementary scalar is
$1/2$ and has rank one; the asserted bound is also one.
The lemma gives the asymmetric claim directly, or it follows by
applying the self-complementary bound to $(M+N)/2$ and using rank
subadditivity.

Taking $k$th roots gives the rate $\sqrt{52}$. To check the strict
comparison without decimal approximations, note that
\[
 52^6=19770609664>17179869184=4^{17},
 \qquad \sqrt2<\frac{17}{12}.
\]
Hence $\sqrt{52}>4^{17/12}>4^{\sqrt2}$, proving the exponent claim.
\end{proof}

This closes a construction family, not the unrestricted square
partial-transpose problem. The algebras with local dimensions two
and three are not covered by this bound, and neither are general
matrices outside the displayed tensor algebra. In particular, the
theorem does not decide whether the square is a price. An exact integer verification of the support certificate is included
in the workbench. Its finite matrix checks verify the input to the
all-copy argument, while the lemma proves the passage to arbitrary
copy numbers.
